\section{Построение суффиксного автомата}

На прошлой лекции мы выяснили, как устроены состояния суффиксного автомата строки $s$: каждое состояние $C$ представляет собой класс контекстно эквивалентных подстрок, состоит из всех суффиксов своей самой длинной строки $\longest(C)$ длины от $\len(\link(C)) + 1$ до $\len(C)$ и однозначно определяется самой строкой $\longest(C)$. При этом подстрока $\alpha$ является строкой $\longest$ некоторого класса тогда и только тогда, когда $\alpha$ "--- префикс $s$, либо $\alpha$ ветвится влево (в $s$ входят $a\alpha$ и $b\alpha$ при $a \ne b$).

Теперь мы готовы построить суффиксный автомат в онлайне, добавляя символы строки $s$ по одному слева направо.

\subsection{Эволюция классов при дописывании символа \texorpdfstring{($s \to sc$)}{(s -> sc)}}

Пусть у нас уже построен суффиксный автомат для строки $s$, и мы дописываем к ней справа один символ $c$. Как изменится множество вершин автомата, то есть множество строк, являющихся $\longest$ в своих классах эквивалентности?

\begin{proposition}[О сохранении старых вершин]
    Любая подстрока $\alpha$ строки $s$, которая была самой длинной в своём классе эквивалентности в строке $s$, остаётся самой длинной в своём классе и в строке $sc$.
\end{proposition}

\begin{proof}
    Проверим критерий из леммы~\ref{lem:longest_crit}:
    \begin{itemize}
        \item если $\alpha$ была префиксом строки $s$, то она, очевидно, остаётся префиксом и более длинной строки $sc$;
        \item если в строку $s$ входили две подстроки $a\alpha$ и $b\alpha$ с различными буквами $a \ne b$, то обе они по-прежнему являются подстроками строки $sc$.
    \end{itemize}
    Следовательно, ни одна старая вершина суффиксного автомата при дописывании буквы $c$ не исчезает и не склеивается с другими вершинами.
\end{proof}

А сколько \emph{новых} строк $\longest$ (то есть новых вершин автомата) может появиться при переходе от $s$ к $sc$?

\begin{lemma}[О числе новых вершин на одном шаге]\label{lem:sa_at_most_2}
    При дописывании одного символа $c$ к строке $s$ множество строк $\longest$ пополняется \textbf{не более чем двумя} строками:
    \begin{enumerate}
        \item всей новой строкой $sc$ (она становится новым префиксом);
        \item возможно, ещё одной строкой $s_0$ "--- самым длинным суффиксом строки $sc$, который уже встречался ранее в строке $s$ как подстрока.
    \end{enumerate}
\end{lemma}

\begin{proof}
    Разобьём все подстроки строки $sc$ на два типа: новые (которых вообще не было в $s$) и старые (которые уже встречались в $s$).

    \textbf{1. Новые подстроки.} Любая подстрока, впервые появившаяся в $sc$, обязана заканчиваться на только что добавленный символ $c$, то есть являться суффиксом строки $sc$. Все такие суффиксы входят в $sc$ ровно один раз "--- на самом правом конце, а значит, имеют один и тот же правый контекст $R = \{\epsilon\}$. Следовательно, все новые подстроки попадают в один общий класс эквивалентности $[sc]$, в котором самой длинной строкой является вся строка $sc$.

    \textbf{2. Старые подстроки.} Пусть $\alpha$ "--- подстрока, которая уже встречалась в $s$, но не была в ней $\longest$, а в строке $sc$ впервые стала $\longest$. Так как $\alpha$ не была префиксом $s$, она не может стать и префиксом $sc$. Значит, в строке $s$ перед всеми вхождениями $\alpha$ всегда стояла одна и та же буква $b$, а в строке $sc$ у $\alpha$ появилось новое вхождение, перед которым стоит другая буква $a \ne b$.

    Но любое новое вхождение подстроки в $sc$ заканчивается в самом правом конце строки $sc$. Значит, строка $\alpha$ обязана быть суффиксом строки $sc$, который уже встречался ранее в $s$.

    Обозначим через $s_0$ \emph{самый длинный} суффикс строки $sc$, который уже встречался в $s$. Тогда любая такая подстрока $\alpha$ является суффиксом строки $s_0$. Покажем, что никакая строка $\alpha$, строго более короткая, чем $s_0$ ($|\alpha| < |s_0|$), не могла впервые стать $\longest$ на этом шаге.

    \begin{center}
    \begin{tikzpicture}[scale=0.95]
        \draw[thick] (0, 0) rectangle (10, 0.7);
        % Old occurrence of s_0
        \draw[thick, fill=NavyBlue!12] (1.2, 0) rectangle (4.2, 0.7);
        \draw[thick, fill=BrickRed!20] (2.2, 0) rectangle (4.2, 0.7) node[pos=0.5, font=\small] {$\alpha$};
        \draw[thick, fill=OliveGreen!25] (1.7, 0) rectangle (2.2, 0.7) node[pos=0.5, font=\small] {$d$};
        \draw[decorate, decoration={brace, amplitude=5pt}, thick] (1.2, 0.85) -- (4.2, 0.85) node[midway, above=5pt, font=\footnotesize] {старое вхождение $s_0$ в $s$};

        % New occurrence of s_0 at the end of sc
        \draw[thick, fill=NavyBlue!12] (7.0, 0) rectangle (10, 0.7);
        \draw[thick, fill=BrickRed!20] (8.0, 0) rectangle (10, 0.7) node[pos=0.5, font=\small] {$\alpha$};
        \draw[thick, fill=OliveGreen!25] (7.5, 0) rectangle (8.0, 0.7) node[pos=0.5, font=\small] {$d$};
        \draw[decorate, decoration={brace, amplitude=5pt}, thick] (7.0, 0.85) -- (10, 0.85) node[midway, above=5pt, font=\footnotesize] {суффикс $s_0$ на конце $sc$};
    \end{tikzpicture}
    \end{center}

    В самом деле, раз $|\alpha| < |s_0|$ и $\alpha$ "--- суффикс $s_0$, то непосредственно перед $\alpha$ внутри строки $s_0$ стоит некоторый символ $d$. Но строка $s_0$ по определению уже встречалась целиком в строке $s$. В том старом вхождении $s_0$ перед суффиксом $\alpha$ стоял точно такой же символ $d$. Значит, новое вхождение $\alpha$ на конце строки $sc$ предваряется тем же самым символом $d$, который уже стоял перед $\alpha$ в строке $s$, и потому не может добавить для $\alpha$ новой предшествующей буквы. Следовательно, единственной старой подстрокой, которая может впервые стать $\longest$, является сама строка $s_0$.
\end{proof}

\subsection{Входящие рёбра и алгоритм перестроения}

Прежде чем выписать алгоритм, сформулируем важное свойство рёбер, входящих в фиксированное состояние автомата.

\begin{lemma}[О структуре входящих переходов]\label{lem:sa_incoming}
    Пусть $q$ "--- любое состояние суффиксного автомата, отличное от начального. Тогда:
    \begin{enumerate}
        \item Все переходы, входящие в состояние $q$, помечены одной и той же буквой $c$ (равной последнему символу любой строки из класса $q$).
        \item Вершины $v$, из которых есть переход $v \xrightarrow{c} q$, образуют непрерывный путь в дереве суффиксных ссылок: $p, \link(p), \dots, p'$.
        \item Если $p$ "--- вершина с наибольшей длиной $\len(p)$ среди всех вершин, из которых есть переход в $q$, то самая длинная строка в состоянии $q$ равна $\longest(q) = \longest(p) + c$, и в частности
        \[
            \len(q) = \len(p) + 1.
        \]
    \end{enumerate}
\end{lemma}

\begin{proof}
    Любая непустая строка в классе $q$ является суффиксом $\longest(q)$, а потому заканчивается на один и тот же последний символ $c$ строки $\longest(q)$. Если отбросить этот последний символ $c$ у всех строк класса $q$, которые имеют длину от $|\shortest(q)|$ до $|\longest(q)|$, мы получим все суффиксы строки $\longest(q)$ без последней буквы $c$, имеющие длины от $|\shortest(q)| - 1$ до $|\longest(q)| - 1$.

    В суффиксном автомате эти суффиксы подряд уложены в цепочку суффиксных ссылок, начиная с некоторой вершины $p$ (содержащей префикс $\longest(q)$ длины $\len(q) - 1$). При этом сама строка $\longest(p)$ при дописывании $c$ тоже обязана попадать в $q$, откуда $\len(p) + 1 = \len(q)$.
\end{proof}

Теперь соберём воедино пошаговый алгоритм добавления символа $c$ к суффиксному автомату строки $s$. Пусть $\text{last}$ "--- состояние, отвечающее всей строке $s$ до добавления $c$.

\textbf{Шаг 1. Создание вершины $\text{cur} = [sc]$ и проведение новых рёбер.}
Создадим новое состояние $\text{cur}$ с длиной $\len(\text{cur}) = \len(\text{last}) + 1$. Нам нужно провести переходы по букве $c$ в $\text{cur}$ из всех суффиксов строки $s$, которые при дописывании $c$ дают новые подстроки (не встречавшиеся в $s$). Все суффиксы строки $s$ лежат на пути суффиксных ссылок $v = \text{last}, \link(\text{last}), \dots, 0$. Будем подниматься по этому пути: пока из текущей вершины $v$ \textbf{нет перехода по букве $c$}, ставим переход $\nxt(v, c) = \text{cur}$ и переходим к $v \la \link(v)$.

\textbf{Шаг 2. Разбор трёх случаев остановки цикла.}
\begin{itemize}
    \item \textbf{Случай 1: цикл дошёл до конца ($v = -1$).} Это значит, что ни из одной вершины на пути до корня (включая сам корень $0$) не было перехода по букве $c$, то есть символ $c$ встретился в строке впервые. В этом случае строка $s_0$ пуста ($s_0 = \epsilon$), все непустые суффиксы $sc$ лежат в $\text{cur}$, поэтому мы полагаем $\link(\text{cur}) = 0$ и завершаем шаг.

    \item Пусть теперь цикл остановился на первой вершине $p$, из которой \textbf{уже существовал} переход $p \xrightarrow{c} q$. Поскольку мы шли по суффиксам строки $s$ в порядке убывания их длины, строка $s_0 = \longest(p) + c$ (длины $\len(p) + 1$) и есть самый длинный суффикс строки $sc$, который уже встречался в $s$. До добавления символа $c$ строка $s_0$ лежала в состоянии $q$.

    \item \textbf{Случай 2: $\len(p) + 1 = \len(q)$.} По лемме~\ref{lem:sa_incoming} это равенство означает, что $s_0 = \longest(q)$, то есть строка $s_0$ \emph{уже являлась} самой длинной строкой в состоянии $q$. Значит, все строки в классе $q$ являются суффиксами $s_0$ (а следовательно, и суффиксами $sc$), поэтому правый контекст всех строк из $q$ синхронно пополняется пустым словом $\epsilon$, и класс $q$ не расщепляется. Достаточно положить $\link(\text{cur}) = q$ и завершить шаг.

    \item \textbf{Случай 3 (клонирование вершины): $\len(p) + 1 < \len(q)$.} В этом случае строка $s_0 = \longest(p) + c$ строго короче, чем $\longest(q)$. Значит, внутри состояния $q$ произошёл раскол:
    \begin{itemize}
        \item более длинные строки класса $q$ (длины от $|s_0| + 1$ до $\len(q)$) \emph{не являются} суффиксами $sc$ "--- их правый контекст не изменился;
        \item более короткие строки класса $q$ (суффиксы $s_0$ длины от $\len(\link(q)) + 1$ до $|s_0|$) \emph{являются} суффиксами $sc$ "--- у них появилось новое вхождение в самом конце строки $sc$.
    \end{itemize}
    Поэтому мы отщепляем от вершины $q$ новую вершину-клон $\text{clone}$, которая заберёт себе все строки из $q$ длины $\le |s_0|$:
    \begin{enumerate}
        \item Полагаем $\len(\text{clone}) = \len(p) + 1$.
        \item Копируем все исходящие переходы из $q$ в $\text{clone}$: $\nxt(\text{clone}) = \nxt(q)$ (так как при дописывании любой буквы справа строки из $\text{clone}$ переходят в те же классы, что и раньше).
        \item Перестраиваем суффиксные ссылки: самая короткая строка в $\text{clone}$ та же, что была в $q$, поэтому $\link(\text{clone}) = \link(q)$; а для оставшейся верхней части $q$ и для новой вершины $\text{cur}$ самым длинным суффиксом вне их классов становится $s_0 = \longest(\text{clone})$, поэтому полагаем $\link(q) = \text{clone}$ и $\link(\text{cur}) = \text{clone}$.
        \item Наконец, перенаправляем входящие переходы: продолжаем идти вверх по суффиксным ссылкам, начиная с вершины $v = p$, и пока $\nxt(v, c) = q$, заменяем переход на $\nxt(v, c) = \text{clone}$. Как только $\nxt(v, c) \ne q$ (или $v = -1$), останавливаемся.
    \end{enumerate}
\end{itemize}

\begin{center}
\begin{tikzpicture}[
    scale=1.0,
    vertex/.style={circle, draw=black!80, thick, minimum size=8.5mm, inner sep=2.5pt, font=\footnotesize},
    clonev/.style={vertex, fill=OliveGreen!20},
    newv/.style={vertex, fill=NavyBlue!20},
    trans/.style={draw=black!80, -Stealth, thick},
    newtrans/.style={draw=OliveGreen!80!black, -Stealth, thick},
    linkedge/.style={draw=BrickRed, -Stealth, dashed, thick}
]
    % Chain of suffix links on the top/left
    \node[vertex] (last) at (-4.5, 1.5) {$s$};
    \node[vertex] (p) at (-1.5, 1.5) {$p$};
    \node[vertex] (lp) at (1.0, 1.5) {$p'$};
    \node[vertex] (root) at (3.5, 1.5) {$0$};

    \draw[linkedge] (last) -- node[above, font=\scriptsize] {$\link$} (p);
    \draw[linkedge] (p) -- node[above, font=\scriptsize] {$\link$} (lp);
    \draw[linkedge] (lp) -- node[above, font=\scriptsize] {$\link$} (root);

    % Bottom vertices: cur, q, clone, link(q)
    \node[newv] (cur) at (-4.5, -1.2) {$\text{cur}$};
    \node[vertex] (q) at (-1.8, -1.2) {$q$};
    \node[clonev] (clone) at (0.8, -1.2) {$\text{clone}$};
    \node[vertex] (lq) at (3.5, -1.2) {$w$};

    % Transitions
    \draw[trans] (last) -- node[left, font=\small\bfseries, text=NavyBlue] {$c$} (cur);
    \draw[trans, gray!50, dotted] (p) -- (q);
    \draw[newtrans] (p) -- node[above right, font=\small\bfseries, text=OliveGreen!80!black] {$c$} (clone);
    \draw[newtrans] (lp) -- node[right, font=\small\bfseries, text=OliveGreen!80!black] {$c$} (clone);

    % Suffix links at bottom
    \draw[linkedge, bend right=30] (cur) to (clone);
    \draw[linkedge] (q) -- (clone);
    \draw[linkedge] (clone) -- (lq);
\end{tikzpicture}
\par\nopagebreak\medskip
\small\textit{Рис. 5.1. Схема клонирования состояния $q$ в Случае 3: переходы по букве $c$ из $p, \link(p), \dots$, ранее ведшие в $q$, перенаправляются в $\text{clone}$, а суффиксные ссылки из $\text{cur}$ и $q$ начинают указывать на $\text{clone}$.}
\end{center}

\begin{lstlisting}[caption={Реализация алгоритма построения суффиксного автомата}]
struct SuffixAutomaton {
    struct State {
        int len = 0;
        int link = -1;
        vector<int> next;

        State(int len = 0, int link = -1) : len(len), link(link), next(SIGMA, -1) {}
    };

    vector<State> st;
    int last;

    SuffixAutomaton() {
        st.emplace_back(0, -1); // корень имеет индекс 0
        last = 0;
    }

    void extend(char ch) {
        int c = ch - 'a';
        int cur = st.size();
        st.emplace_back(st[last].len + 1, -1);

        int p = last;
        while (p != -1 && st[p].next[c] == -1) {
            st[p].next[c] = cur;
            p = st[p].link;
        }

        if (p == -1) {
            // Случай 1: символ c встретился впервые
            st[cur].link = 0;
        } else {
            int q = st[p].next[c];
            if (st[p].len + 1 == st[q].len) {
                // Случай 2: $s_0$ уже является самой длинной строкой в $q$
                st[cur].link = q;
            } else {
                // Случай 3: расщепление состояния q и создание клона
                int clone = st.size();
                st.emplace_back(st[p].len + 1, st[q].link);
                st[clone].next = st[q].next; // копируем исходящие переходы

                while (p != -1 && st[p].next[c] == q) {
                    st[p].next[c] = clone;
                    p = st[p].link;
                }
                st[q].link = st[cur].link = clone;
            }
        }
        last = cur;
    }
};
\end{lstlisting}

\subsection{Оценка числа вершин и рёбер суффиксного автомата}

Сколько памяти занимает построенный суффиксный автомат? Докажем, что несмотря на то что в строке длины $n$ бывает $\Theta(n^2)$ различных подстрок, и число вершин, и число переходов в суффиксном автомате линейны по $n$.

\begin{theorem}[О размере суффиксного автомата]\label{thm:sa_size}
    Для любой строки $s$ длины $n \ge 2$ суффиксный автомат содержит не более $2n - 1$ вершин и не более $3n - 4$ переходов (грубо: $|V| \le 2n$ и $|E| \le 3n$).
\end{theorem}

\begin{proof}
    \textbf{1. Оценка числа вершин.} До добавления символов в автомате есть $1$ вершина (корень $0$). При добавлении первого символа $s[0]$ мы всегда попадаем в Случай 1 и создаём ровно $1$ вершину. При добавлении второго символа $s[1]$ клонирование также невозможно (так как у единственной нетривиальной вершины длина равна $1$, а $\len(0) + 1 = 1$), поэтому снова создаётся ровно $1$ вершина. На каждом из оставшихся $n - 2$ шагов по лемме~\ref{lem:sa_at_most_2} создаётся не более $2$ вершин ($\text{cur}$ и, возможно, $\text{clone}$). Итого общее число вершин не превосходит
    \[
        |V| \le 1 + 1 + 1 + 2(n - 2) = 2n - 1.
    \]

    \textbf{2. Оценка числа рёбер.} Заметим, что при переходе по любому ребру $u \xrightarrow{c} v$ длина самой длинной строки строго возрастает ($\len(v) \ge \len(u) + 1$), поэтому суффиксный автомат является ориентированным ациклическим графом (DAG), в котором все вершины достижимы из корня $0$.

    Выделим в суффиксном автомате произвольное ориентированное остовное дерево $T$ с корнем в начальной вершине $0$. В дереве $T$ ровно $|V| - 1 \le 2n - 2$ рёбер (назовём их \textbf{остовными}).

    \begin{center}
    \begin{tikzpicture}[
        scale=1.0,
        vertex/.style={circle, draw=black!80, thick, minimum size=6.5mm, inner sep=1pt, font=\footnotesize},
        term/.style={vertex, fill=NavyBlue!20, double, double distance=1.2pt},
        treeedge/.style={draw=black!80, -Stealth, very thick},
        nonedge/.style={draw=BrickRed, -Stealth, thick, dashed},
        anypath/.style={draw=OliveGreen!80!black, -Stealth, thick, decorate, decoration={snake, amplitude=1pt, segment length=8pt, post length=3pt}}
    ]
        \node[vertex] (0) at (0, 0) {$0$};
        \node[vertex] (u) at (3.4, 0.55) {$u$};
        \node[vertex] (v) at (6.4, 0.55) {$v$};
        \node[term]   (t) at (9.8, 0) {$t$};

        \draw[treeedge] (0) to[bend left=16] node[pos=0.42, above=4pt, font=\footnotesize] {путь в $T$ ($\alpha_e$)} (u);
        \draw[nonedge] (u) -- node[above=3pt, font=\footnotesize\bfseries, text=BrickRed] {$e = (u \xrightarrow{c} v)$} (v);
        \draw[anypath] (v) to[bend left=12] node[pos=0.58, above=4pt, font=\footnotesize, text=OliveGreen!80!black] {любой путь ($\beta_e$)} (t);
    \end{tikzpicture}
    \end{center}

    Оценим теперь количество \textbf{неостовных рёбер}. Каждому неостовному ребру $e = (u \xrightarrow{c} v)$ сопоставим непустой суффикс $w_e$ строки $s$, составленный из трёх участков пути:
    \begin{enumerate}
        \item сначала пройдём от корня $0$ до вершины $u$ по единственному пути в остовном дереве $T$ (пусть на этом пути написана строка $\alpha_e$);
        \item затем шагнём по самому неостовному ребру $e = (u \xrightarrow{c} v)$ (дописывая символ $c$);
        \item наконец, из вершины $v$ дойдём по любому ориентированному пути в автомате до какой-нибудь терминальной вершины $t$ (пусть на этом участке написана строка $\beta_e$).
    \end{enumerate}
    Поскольку весь маршрут ведёт из корня $0$ в терминальную вершину $t$, прочитанная на нём строка $w_e = \alpha_e c \beta_e$ допускается автоматом, а значит, является непустым суффиксом строки $s$.

    Докажем, что разным неостовным рёбрам $e_1 \ne e_2$ соответствуют \textbf{попарно различные} суффиксы $w_{e_1} \ne w_{e_2}$. В самом деле, пусть нам известен суффикс $w_e$. Так как суффиксный автомат детерминирован, любой строке $w_e$ соответствует \emph{единственный} путь в автомате, выходящий из корня $0$. Пройдём по этому пути от корня: сначала мы идём по рёбрам остовного дерева $T$ (участок $\alpha_e$), и \textbf{первым же ребром на этом пути, не принадлежащим остовному дереву $T$, будет в точности ребро $e$}! Значит, по суффиксу $w_e$ ребро $e$ восстанавливается однозначно, и отображение $e \mapsto w_e$ инъективно.

    Но всего у строки $s$ ровно $n$ непустых суффиксов (причём самый длинный суффикс $s$ целиком обязан иметь хотя бы один путь без циклов; более того, хотя бы один суффикс достигается чисто по остовному дереву). Следовательно, неостовных рёбер не более $n$, а общее число всех рёбер удовлетворяет оценке $|E| \le (2n - 2) + n < 3n$.
\end{proof}
